\documentclass[11pt,a4paper]{article}
\usepackage[utf8]{inputenc}
\usepackage[vietnamese]{babel}
\usepackage{amsmath,amssymb}
\usepackage{geometry}
\geometry{a4paper, total={180mm,260mm}, left=15mm, right=15mm, top=18mm, bottom=18mm}

\usepackage{tikz}
\usetikzlibrary{shapes.geometric, arrows.meta, positioning, calc}
\tikzset{
  arr/.style={-{Stealth[scale=0.8]}},
  every node/.append style={align=center}
}
\usepackage{tcolorbox}
\tcbuselibrary{skins,breakable}
\usepackage{fontawesome}
\usepackage{fancyhdr}
\usepackage{enumitem}

\pagestyle{fancy}
\fancyhf{}
\renewcommand{\headrulewidth}{0.5pt}
\renewcommand{\footrulewidth}{0.5pt}
\fancyhead[L]{\small\bfseries\color{violet!85!black}\faBook\ SỔ TAY KIẾN THỨC TOÁN 2}
\fancyhead[R]{\small\bfseries\color{violet!85!black}Lớp Toán Thầy Lĩnh -- loptoanthaylinh.vn}
\fancyfoot[C]{\small\color{gray}-- Sổ tay kiến thức: Bài 46. Khối trụ, khối cầu --}

\setlength{\parindent}{0pt}
\setlength{\parskip}{0pt}

\begin{document}

% --- BANNER BÀI HỌC ---
\begin{tcolorbox}[
  enhanced,
  colback=violet!8!white,
  colframe=violet!70!black,
  arc=3mm, boxrule=1.2pt,
  left=3mm, right=3mm, top=2mm, bottom=2mm,
  halign=center
]
  {\LARGE\bfseries\color{violet!90!black}\faBookmark\ BÀI 46: KHỐI TRỤ, KHỐI CẦU}\\[1mm]
  {\itshape\footnotesize\color{black!75} Nhận biết hình dạng, đặc điểm và vật thật thực tế của khối trụ và khối cầu trên 1 trang A4}
\end{tcolorbox}

\vspace{-2mm}

\noindent
\begin{minipage}[t]{0.49\textwidth}\vspace{0pt}
\begin{tcolorbox}[
  enhanced,
  colback=violet!4!white,
  colframe=violet!75!black,
  arc=2mm, boxrule=0.8pt,
  left=2mm, right=2mm, top=1.8mm, bottom=1.8mm,
  title={\footnotesize\bfseries\color{white}\faKey\ 1. Khối trụ}
]
{\footnotesize
\begin{itemize}\setlength{\itemsep}{1pt}\setlength{\leftskip}{-2.5mm}
  \item \textbf{Đặc điểm cấu tạo:}\\
  - Có \textbf{hai mặt đáy} là hai hình tròn bằng nhau phẳng.\\
  - Có một \textbf{mặt xung quanh} cong tròn.
  \item \textbf{Khả năng chuyển động:}\\
  - Có thể \textbf{trượt} (khi đặt trên mặt đáy phẳng).\\
  - Có thể \textbf{lăn} (khi đặt nằm ngang trên mặt cong).\\
  - Có thể \textbf{chồng lên nhau} được (đặt đáy lên đáy).
  \item \textbf{Vật thật:}\\
  Hộp sữa đặc, lon nước ngọt, khúc gỗ tròn, cục pin, trống trường,...
\end{itemize}
}
\end{tcolorbox}
\end{minipage}
\hfill
\begin{minipage}[t]{0.49\textwidth}\vspace{0pt}
\begin{tcolorbox}[
  enhanced,
  colback=teal!4!white,
  colframe=teal!75!black,
  arc=2mm, boxrule=0.8pt,
  left=2mm, right=2mm, top=1.8mm, bottom=1.8mm,
  title={\footnotesize\bfseries\color{white}\faSliders\ 2. Khối cầu}
]
{\footnotesize
\begin{itemize}\setlength{\itemsep}{1pt}\setlength{\leftskip}{-2.5mm}
  \item \textbf{Đặc điểm cấu tạo:}\\
  - Toàn bộ bề mặt là một \textbf{mặt cong khép kín hoàn hảo}.\\
  - Không có đỉnh, không có cạnh và không có mặt phẳng đáy nào.
  \item \textbf{Khả năng chuyển động:}\\
  - Có thể \textbf{lăn về mọi hướng}.\\
  - \textbf{Không thể chồng} đứng lên nhau được (sẽ bị trượt đổ).
  \item \textbf{Vật thật:}\\
  Quả bóng đá, viên bi thủy tinh, quả địa cầu, viên ngọc trai, quả cam,...
\end{itemize}
}
\end{tcolorbox}
\end{minipage}

\vspace{-1.5mm}
\begin{tcolorbox}[
  enhanced,
  colback=orange!4!white,
  colframe=orange!80!black,
  arc=2mm, boxrule=0.8pt,
  left=2mm, right=2mm, top=1.5mm, bottom=1.5mm,
  title={\footnotesize\bfseries\color{white}\faSitemap\ 3. Bảng đối chiếu so sánh Khối trụ và Khối cầu}
]
\centering
\begin{tikzpicture}[x=1.5cm, y=0.8cm]
  \node[rectangle, rounded corners=2pt, draw=violet!80!black, fill=violet!15, text width=4.5cm, align=center, font=\scriptsize\bfseries] at (0,0) {
    \textbf{KHỐI TRỤ}\\$\bullet$ 2 đáy phẳng hình tròn.\\$\bullet$ Vừa lăn được, vừa trượt được.\\$\bullet$ Chồng lên nhau được.
  };
  \node[rectangle, rounded corners=2pt, draw=teal!80!black, fill=teal!15, text width=4.5cm, align=center, font=\scriptsize\bfseries] at (4.8,0) {
    \textbf{KHỐI CẦU}\\$\bullet$ Mặt cong toàn bộ.\\$\bullet$ Lăn được về mọi hướng.\\$\bullet$ Không chồng lên nhau được.
  };
\end{tikzpicture}
\end{tcolorbox}

\vspace{-1.5mm}
\noindent
\begin{minipage}[t]{0.49\textwidth}\vspace{0pt}
\begin{tcolorbox}[
  enhanced,
  colback=green!4!white,
  colframe=green!60!black,
  arc=2mm, boxrule=0.8pt,
  left=2mm, right=2mm, top=1.8mm, bottom=1.8mm,
  title={\footnotesize\bfseries\color{white}\faCheckCircle\ 4. Ví dụ thực hành chuẩn mực}
]
{\footnotesize
\textbf{Phân loại đồ vật:}\\
Xếp các vật sau vào nhóm thích hợp:\\
\textit{Lon cocacola, quả bóng tennis, hộp chè tròn, viên bi.}\\
- \textbf{Khối trụ:} Lon cocacola, hộp chè tròn.\\
- \textbf{Khối cầu:} Quả bóng tennis, viên bi.
}
\end{tcolorbox}
\end{minipage}
\hfill
\begin{minipage}[t]{0.49\textwidth}\vspace{0pt}
\begin{tcolorbox}[
  enhanced,
  colback=red!4!white,
  colframe=red!75!black,
  arc=2mm, boxrule=0.8pt,
  left=2mm, right=2mm, top=1.8mm, bottom=1.8mm,
  title={\footnotesize\bfseries\color{white}\faExclamationTriangle\ 5. Mẹo tránh bẫy trắc nghiệm}
]
{\footnotesize
\begin{itemize}\setlength{\itemsep}{1pt}\setlength{\leftskip}{-2.5mm}
  \item \textbf{Phân biệt Khối cầu và Hình tròn:}\\
  - Hình tròn: Hình phẳng vẽ trên giấy 2D.\\
  - Khối cầu: Vật thể không gian 3D cầm nắm được.
  \item Khối trụ chỉ lăn được khi đặt nằm nghiêng, đặt đứng thì chỉ trượt được.
\end{itemize}
}
\end{tcolorbox}
\end{minipage}


\end{document}
